\subsection{A rank-four nonabelian obstruction and its compatible replacement}

\begin{statement}{D.30}{Proved}{The first new rank-four layer and an actual class-three kernel}
The quotient map
\[
 Q\longrightarrow Q_0:=Q/O_2(Q)
       \simeq S_3\times C_3                              \tag{D.30.1}
\]
has kernel $O_2(Q)\simeq C_2^2\times Q_8$ of order 32.
Put $K=\ker(B\to Q_0)$. Then $K$ is free of rank 4,
$N\triangleleft K$ has index 32, and the relative pro-two
completion $\mathcal P_2$ has free pro-two kernel $K^{(2)}$
over $Q_0$. Its index 32 subgroup is the rank 97 group
$N^{(2)}$ used previously.

Write $D_j(K)$ for the mod-two Zassenhaus filtration and set
\[
 T_0=B/D_6(K),\qquad T_1=B/D_7(K).
                                                               \tag{D.30.2}
\]
Their orders are $576\cdot2^{303}$ and
$576\cdot2^{993}$, respectively. The new central layer
$Z=D_6(K)/D_7(K)$ has dimension 690. The image
\[
 W=\gamma_3N\,D_7(K)/D_7(K)\subset Z                  \tag{D.30.3}
\]
has dimension 240. Central cusp corrections have rank 227
on $W$ and 648 on $Z$.

Both the specific alternate $H_2$ representative recorded
in D.27 and the specific changed $H^{(1)}$ representative
retained in D.29 have no genuine full profinite lift.
This conclusion uses their retained $G_6$ action jointly
with their $T_0$ action. It follows after exhausting the
entire lower $N$-centralizer freedom: the image of $N$ in
$T_0$ has central Frattini subgroup of dimension 225 and
Frattini quotient of dimension 78; its full $s$-centralizer
is generated by 152 explicitly computed elements. Their
defects raise the $Z$ correction rank to 654, and each
recorded cusp target remains outside that image.
Separating linear functionals are saved with the exact
rank certificates. This excludes these representatives,
not every action with their CM pair or their $G_6$ data.

There is a further explicit replacement with the same
homology masks, right-centralizer data, $G_3$ conjugators
and extra CM pair. Its simultaneous correction matrix
has rank 28: fifteen conditions retain the $\mathbb Z/4$
class-two lift and thirteen remove the new $W$ obstruction.
The resulting 2944468-letter cusp word vanishes both in
the full rank 97 $\mathbb Z/4$ model and in the full
rank-four Magnus model through degree six. Its multiplier
is 1. Thus it gives compatible automorphisms on
$H^{(2)}$ and $T_1$, with exact common quotient
\[
 C_*=B/(\gamma_3N\,D_7(K)),\qquad
 |C_*|=576\cdot2^{753},
 \quad |H^{(2)}\times_{C_*}T_1|=576\cdot2^{9843}.
                                                               \tag{D.30.4}
\]

The replacement extends through the entire stable tower
of D.29 while retaining its $T_1$ action and conjugators.
Its two-primary target is the precise quotient
\[
 \mathcal S_2=
 \mathcal P_2/\bigl(
 \overline{\gamma_3N^{(2)}}\cap D_7(K^{(2)})\bigr)
 \simeq\mathcal P_2^{[2]}\times_{C_*}T_1.              \tag{D.30.5}
\]
Its kernel over $Q$ is nilpotent of class exactly three:
its third lower-central subgroup is $W\simeq C_2^{240}$,
and its fourth is trivial. The same action extends over
every odd relative pro-prime kernel and the compatible
mixed-prime quotient $J$ of D.28. One global cusp gauge
gives nonarithmetical finite shadows on the characteristic
refinements retaining the $G_6$ and CM data. The genuine
full profinite realization of this replacement and of
the extra CM pair remains undecided.
\end{statement}
\begin{proof}
The map (D.30.1) is induced by reduction modulo two and
by reduction modulo three followed by
$\mathrm{SL}_2(\mathbb F_3)\to
\mathrm{SL}_2(\mathbb F_3)/Q_8=C_3$.
Both maps kill the global sign, and Chinese remainders
make their joint image $S_3\times C_3$. Its kernel has
order 32 and is exactly the subgroup used in D.27.
The images of $s$ and $r$ retain orders two and three.
Thus $K$ acts freely on the Bass--Serre tree. Its
quotient graph has 18 edges and 15 vertices, so its
fundamental group has rank 4. The same graph argument
gives $[K:N]=32$ and rank 97 for $N$.

The induced pro-two topology on $N$ is its full pro-two
topology. Indeed the intersection of the finitely many
$K$-conjugates of a normal subgroup of two-power index
in $N$ still has two-power index in $N$, and then in $K$.
Conversely restriction of a finite two-group quotient
of $K$ gives one of $N$. Taking the finite normal cores
under $B$ preserves this property. This proves the
relative-completion comparison and its specified kernel
maps, without identifying completions from ranks alone.

The Zassenhaus quotient is evaluated by the Magnus map
over $\mathbb F_2$, retaining all noncommutative monomials
of degree less than its index. The exact Hall-power
basis is \cite[Proposition 4.4 and its proof]{dEfrat2023};
its Lyndon standard factorization is Lemma 2.2 there.
For four free generators the dimensions in degrees
one through six are
\[
                    4,10,20,70,204,690.                \tag{D.30.6}
\]
For degree six the program uses the 670 ordinary
Lyndon brackets and the squares of the 20 degree-three
Lyndon brackets. It checks their independence in the
4096-dimensional associative tensor space. Summing
(D.30.6) gives308 through degree five and 998 through
degree six. Dividing the corresponding kernel orders
by $[K:N]=32$ proves the orders in (D.30.2).

The map $K/D_2(K)\to O_2(Q)^{\mathrm{ab}}$ is an
isomorphism of four-dimensional spaces, so
$N\subset D_2(K)$. The group $O_2(Q)$ has exponent four
and central commutator subgroup of order two, hence
$D_3(K)\subset N$. Its degree-two relation space
$R_2\subset D_2(K)/D_3(K)$ has dimension nine.
All these maps are also checked from the actual 97
words of the original $N$ basis in the new rank-four
graph. Their degree-one terms vanish and their
degree-two terms span exactly $R_2$.
The filtration laws give
\[
 \gamma_3N\subset D_6(K),\qquad
 [N,N]^2N^4\subset D_8(K).                            \tag{D.30.7}
\]
These laws follow directly from multiplication of
the Magnus series: commutator degrees add and
squaring doubles the least degree. The degree-six
image of $\gamma_3N$ is therefore
$[R_2,[R_2,R_2]]$. The code forms all 729 ordered triple
brackets and finds rank 240. It retains a basis made
of actual triple commutators of $N$ words. This proves
(D.30.3) and provides a quotient map $H^{(1)}\to T_0$.

The initial selected cusp word is evaluated in the
rank-four model. Its terms below degree six vanish,
and its degree-six term lies in $W$, as (D.30.7)
requires. For central corrections $b,a\in Z$ the
cusp-error map is
\[
                 (1-s)b+s(1-r)a.                     \tag{D.30.8}
\]
Here conjugation on $Z$ factors through $Q_0$; the
kernel $K$ acts trivially. Tensor substitution of the
exact four-dimensional graph matrices gives ranks 227
on $W$ and 648 on $Z$. Both selected targets fail these
fixed-conjugator tests.

We next exhaust the remaining lower conjugators, so
the negative conclusion has a precise source. Let
$P$ be the image of $N$ in $T_0$. It is class two of
exponent four by (D.30.7), and has order $2^{303}$.
Its Frattini subgroup
$A_0=P^2[P,P]$ is central elementary abelian. The
program evaluates all 97 generator squares and 4656
pairwise commutators in the rank-four Magnus model.
Their exact span has dimension 225. Collection modulo
that span gives the 78-dimensional elementary abelian
quotient $V_0=P/A_0$; the dimensions sum to 303, and
the constructed words generate all of $P$.

Write $s$ also for its conjugation action. The connecting
map
\[
 V_0^s\longrightarrow A_0/(1-s)A_0,
 \qquad \bar c\longmapsto[c,s]\bmod(1-s)A_0           \tag{D.30.9}
\]
is well-defined. It is additive because $A_0$ is central
in $P$, and changing a lift changes the defect by
$(1-s)A_0$. Exact computation gives
\[
 \dim A_0^s=117,\quad \dim V_0^s=42,\quad
 \operatorname{rank}(D.30.9)=7.                       \tag{D.30.10}
\]
Choose a basis of its 35-dimensional kernel and correct
each lift by an element of $A_0$ to centralize $s$.
Together with a basis of $A_0^s$, these are 152 actual
generators of $C_P(s)$: their images generate the
quotient and their central kernel is precisely $A_0^s$.
The code computes every corrected generator with full
degree-six terms, not just its image in $P$.

Its further defect in $Z$ is additive modulo
$(1-s)Z$, since $N\subset K$ acts trivially on $Z$.
Thus adjoining these 152 columns accounts for the
entire lower $s$-centralizer freedom. The other
elliptic centralizer lifts through the two-kernel by
coprime order-three conjugacy. Moreover
$q_0(A)=q_0(\beta)=1$ for the specified CM pair, so
no omitted conjugation changes the new defect columns.
The full augmented matrix has rank 654. Each of the
two selected target vectors is outside its image;
the saved separating functional annihilates every
column and takes value one on that target.

There is no remaining multiplier column. The actual
rank-four word $(sr)^6$ has degree-one vector
$(1,1,0,1)$. Its fourth power has a nonzero degree-four
term, while its eighth power has no term through
degree six. Thus $t=sr$ has order 48 in both $T_0$
and $T_1$. Every multiplier change preserving the
selected $T_0$ action is already trivial on $T_1$.

For a genuine full lift, the original Herfort--Ribes
torsion theorem used in D.27 supplies elliptic
conjugators in $\widehat F$. The retained $G_6$
action forces their $Q$ coordinates by (D.25.6).
This use of $G_6$ is essential: $T_0$ alone does not
retain all its homology data. Differences between
the chosen and new conjugators therefore lie in $N$,
and the preceding finite centralizer computation
exhausts their reductions. Undoing the specified cusp
gauge retains the cusp power. The two nonzero
separating values prove the stated absence of a
genuine lift for these selected representatives.
No assertion that arbitrary torsion in $T_1$ is
conjugate into its elliptic factors is used.

We now construct the replacement. It changes central
coordinates explicitly, retaining the homology data.
Use $C_1\simeq(\mathbb Z/4)^{4753}$ from D.29, with
basis the squares and commutators of the 97 $N$
generators. Their images in $K/D_7(K)$ have degree
at least four, so commute there as well. Hence their
cusp effects are linear in both target groups.
The degree-six action on these images is computed
from the actual elliptic conjugates, including their
degree-one and degree-two path terms; replacing those
conjugates by their $Q_0$ matrices would lose terms.
For the $A$-column one may use the action of $t$ in
place of the product of the two elliptic conjugates:
their discrepancy has degree six, whose conjugation
is trivial on elements starting in degree four.

The 4769-dimensional relation space of the mod-two
central correction matrix has two exact obstruction
maps. The first is D.29's fifteen-dimensional
obstruction to lifting that correction modulo four.
The second is the thirteen-dimensional quotient of
$W$ by (D.30.8). The joint map has rank 28. The code
solves for first obstruction zero and second
obstruction equal to the selected cusp defect.
An upper central correction then preserves the full
$\mathbb Z/4$ cusp relation. Its image starts in
$D_8(K)$, so it does not affect the degree-six test.
Finally the retained triple-commutator basis of $W$
removes that degree-six error. Those corrections
vanish identically in every class-two quotient of $N$.
This gives one pair of actual words satisfying both
equations simultaneously.

The result is verified directly by
\texttt{d\_cm\_rank4\_joint.py}. Its 2944468-letter
$B$ word collapses to 562814 letters in the rank-four
free basis. Every degree-one and degree-two coordinate
vanishes in the rank 97 $\mathbb Z/4$ model, and every
coordinate in degrees one through six vanishes in
the rank-four mod-two model. The multiplier is 1.
The map on $N^{\rm ab}/2$ retains rank 97. The map on
$K/D_2(K)$ is the identity, because its identification
with $O_2(Q)^{\rm ab}$ is preserved and the action on
$Q$ is the identity. The finite two-kernel Frattini
argument therefore proves automorphisms of both quotients.

The kernel defining $H^{(2)}$ is
$\gamma_3N[N,N]^4N^8$. The last two factors lie in
$D_{16}(K)$, hence in $D_7(K)$. Therefore its product
with $D_7(K)$ is exactly $\gamma_3N D_7(K)$.
This proves the common-quotient formula (D.30.4).
The kernel of $T_1\to C_*$ is exactly $W$ of order
$2^{240}$, giving all three stated finite orders.
The common conjugator words prove compatibility on
$C_*$, and the normal-subgroup product formula proves
surjectivity onto the fibre product.

Restart D.29's stable iteration from this new seed.
Every later correction in $C_k$, $k\geq2$, lies in
$D_8(K)$: its generators are powers $N^{2^k}$ and
$[N,N]^{2^{k-1}}$. The corresponding multiplier
change acts trivially on $T_1$. Thus the entire
iteration retains this finite action and its exact
conjugators. The inverse-limit normal-subgroup
formula gives (D.30.5). Its third lower-central
subgroup over $Q$ is the image of $\gamma_3N$,
namely $W$. Since $[N,\gamma_3N]\subset D_8(K)$,
its fourth lower-central subgroup is trivial.
The nonzero dimension 240 proves that the class is
exactly three.

The odd-prime centralizer construction of D.29 is
unchanged. The mixed quotient $J$ still agrees on
$H\times_QG_3$: all changes lie in $N^2[N,N]$ or
$\gamma_3N$ on the two-primary side, and the
three-primary conjugators are retained as separately
chosen components of the fibre product. The displayed
discrete two-primary correction word need not itself
have the original $G_3$ abelianization. The maximal
common quotient of $J$ with the new primewise system is
the same one, because its finite kernel over $Q$
has maximal nilpotent quotient computed in D.28.
The combined conjugators therefore lift to a single
pair in the resulting fibre product. The one global
cusp gauge proves the stated compatible finite
shadow family and its nonarithmeticality on the
refinements retaining the CM data. Further finite
quotients remain in D.22's genuine-lift test.
\end{proof}
